%
% File: conclusion.tex
%
\chapter{Conclusion}
\label{chap:conclusion}
\setlength{\parskip}{1em}

\section{Summary of Findings}

This work presented a comprehensive comparative study of recommender system approaches for e-commerce, specifically focusing on the Headphones subcategory of Amazon Electronics reviews.

Our main findings are:
\begin{enumerate}
    \item \textbf{Neural methods show promise}: NeuMF (NCF) achieved the highest ranking performance (NDCG@10 = 0.0059), outperforming both classical Matrix Factorization and baseline approaches.
    
    \item \textbf{Sparsity remains challenging}: With 97.4\% sparsity, traditional collaborative filtering methods like Item-KNN and MF struggled to learn meaningful patterns.
    
    \item \textbf{Popularity is a strong baseline}: The simple popularity-based recommender achieved competitive results (NDCG@10 = 0.0050), highlighting the importance of the long-tail distribution in e-commerce data.
    
    \item \textbf{Cold-start is significant}: Approximately 50\% of test users and items were not seen during training, underscoring the importance of cold-start handling in production systems.
\end{enumerate}

\section{Limitations}

This study has several limitations:
\begin{itemize}
    \item The dataset size after filtering is relatively small (2,695 interactions), which may limit the generalizability of findings
    \item Hybrid models incorporating content features (e.g., LightFM) were not fully evaluated due to technical constraints
    \item The study focuses on offline evaluation; online A/B testing would provide stronger evidence of practical utility
\end{itemize}

\section{Future Work}

Future research directions include:
\begin{enumerate}
    \item \textbf{Hybrid approaches}: Incorporating item metadata (category, brand, price) to address cold-start
    \item \textbf{Sequential models}: Exploring transformer-based architectures (e.g., SASRec, BERT4Rec) for session-based recommendation
    \item \textbf{Larger datasets}: Evaluating on larger subcategories or the full Amazon dataset
    \item \textbf{Explainability}: Investigating interpretable recommendation methods
\end{enumerate}

\section{Reproducibility}

All code, data processing scripts, and trained models are available in the accompanying repository. The experiments can be reproduced using the provided \texttt{run\_all.sh} script and notebook files.